\addchap{Executive summary}\label{ch:summary}

\paragraph{Research motivation.}
An experimental study of crude dichloromethane (DCM) seed and leaf extracts of
\textit{Moringa oleifera} Lam.\ observed that the extracts changed how irradiated cells
proliferated, and that the change differed between an apparently normal rodent fibroblast
line (V79) and two human carcinoma lines (DU145 and MCF-7). A three-line experiment cannot
say whether such differences reflect the cells' \emph{intrinsic} radiation response, nor
whether measurable molecular features explain it. This project asked that broader question
on a defensible endpoint and with public data, and then returned to the plant to ask which
of its documented constituents deserve experimental testing.

\paragraph{Computational question.}
Does baseline gene expression, summarised as 50 Hallmark pathway scores, predict
experimentally measured clonogenic radiosensitivity across cancer cell lines better than
prespecified simple baselines, once related cell lines are kept together and tissue of
origin is already taken into account?

\paragraph{Data and design.}
We assembled 241 eligible cell-line records, forming 237 independent biological components,
by linking literature-derived clonogenic surviving fractions at 2~Gy (CLARIFi mSF2) to
baseline expression (DepMap 24Q4). Cell-line identity and relatedness were resolved through
Cellosaurus, and pathway features were computed from the MSigDB Hallmark gene sets. Five
prespecified models, from a global mean to tissue, proliferation and expression together,
were compared under ten repeats of five-fold cross-validation grouped by biological
component and stratified by tissue, with every tuning and scaling step confined to training
folds. All 250 model-fold tasks completed.

\paragraph{Principal findings.}
Adding expression features to a tissue-only baseline reduced the mean absolute error (MAE)
by \textbf{0.00383 mSF2 units, 2.85\% of the baseline error}, and the richer model had the
lower error in \textbf{nine of ten} cross-validation repeats. Out-of-sample $R^2$ stayed
close to zero for every model ($-0.018$ to $+0.013$). Expression therefore carries a small
amount of usable information beyond tissue of origin, while overall predictive power is
weak --- a result that agrees with the strongest published benchmark in the field.
Superiority was not formally tested: a protocol deviation, recorded before any result was
seen, deferred the coverage simulation needed to validate an interval method, and no
negative-control draws were run. The improvement is reported as a descriptive estimate.

\paragraph{Robustness and generalisation.}
Three further analyses of the same frozen cohort tested how durable this small advantage
is. A random forest (M3) did \textbf{worse} than the linear offset model (gain over the
tissue baseline 0.00275 against 0.00383). On transformed outcome scales the advantage kept
its direction but almost vanished on the neglog scale (0.00023 in neglog units) and was
larger in absolute terms on the logit scale (0.00980 in logit units), so its size depends on
how the outcome is expressed. When one tissue at a time was held out, an expression-only model
reached out-of-sample $R^2$ of \textbf{+0.031} against a global mean, the only clearly
positive $R^2$ in the project; it nevertheless lost to the global mean when Pancreas
($-0.0142$) or Esophagus/Stomach ($-0.0030$) was held out, so transfer to an unseen tissue
is conditional on the tissue.

\paragraph{Explainability.}
Held-out permutation importance ranked the 50 Hallmark scores. Seven reached the
prespecified exploratory tier~C, led by the p53 pathway; the higher tiers were unreachable
by design, because they require a formally demonstrated prediction gain and at least 1,000
negative-control draws. Importance magnitudes are as small as the predictive gain itself,
rankings are only moderately stable across fits, and the DNA repair, G2M checkpoint and
hypoxia sets that radiobiology would lead one to expect were not prioritised.

\paragraph{Compound prioritisation.}
\textit{M.~oleifera} seed and leaf constituents were curated from primary analytical
studies, with the extraction solvent of every report recorded, and linked to
radiation-response biology through a four-dimension score of 0--8 points fixed before any
score was computed. Eighteen compounds were ranked and five reached tier~A. The curation
exposed the central mismatch of the project: \textbf{niazimicin is the only constituent
with a structurally confirmed identification in a dichloromethane fraction of both seed
and leaf, yet no study has tested the purified compound in any radiation-response process;
conversely, the compounds with the strongest radiation evidence have never been found in a
dichloromethane extract.} Two prespecified safeguards visibly bound: chlorogenic acid scored
7 of 8 but is held at tier~C because one study reports that it protects \emph{tumour} cells
from radiation, and benzyl isothiocyanate is excluded despite strong radiation evidence
because it is a root, not a seed or leaf, constituent. A tier summarises the state of the
published evidence; it never asserts radioprotection.

\paragraph{Significance.}
The project delivers a reproducible, auditable pipeline, an honest benchmark whose
negative-leaning result agrees with the best prior work, and a traceable route from the
original radiobiology to a short list of experiments that could settle the question. Its
value lies in the rigour of its methods and the transparency of its evidence rather than in
a strong performance claim.

\begin{keybox}[title={Recommended next steps.}]
Characterise the original DCM extracts analytically (niazimicin and moringin first);
complete the deferred coverage simulation and the negative controls; and test
$\alpha$-tocopherol and $\beta$-sitosterol, purified and against whole extract, on
clonogenic survival in matched normal and tumour cells over a dose range.
\end{keybox}
